% 작업장 공통 서식 (W2 이후). tectonic으로 컴파일.
\documentclass[11pt,a4paper]{article}
\usepackage{fontspec}
\usepackage{xeCJK}
\setCJKmainfont{AppleMyungjo}
\setCJKsansfont{Apple SD Gothic Neo}
\setCJKmonofont{Apple SD Gothic Neo}
\xeCJKsetup{CJKspace=true}
\usepackage{amsmath,amssymb,amsthm,mathtools}
\usepackage[margin=23mm]{geometry}
\usepackage{booktabs,tabularx,array}
\usepackage{enumitem}
\usepackage{xcolor}
\usepackage[most]{tcolorbox}
\usepackage[hidelinks]{hyperref}
\linespread{1.25}
\setlength{\parindent}{0pt}
\setlength{\parskip}{0.55em}

\definecolor{ink}{HTML}{1B2433}
\definecolor{acc}{HTML}{2C55C9}
\definecolor{hot}{HTML}{B8650A}
\newtcolorbox{thmbox}{colback=acc!4,colframe=acc!60,boxrule=0.5pt,arc=1pt,left=6pt,right=6pt,breakable}
\newtcolorbox{bridge}{colback=hot!5,colframe=hot!70,boxrule=0.5pt,arc=1pt,left=6pt,right=6pt,breakable}
\newtcolorbox{openprob}{colback=ink!3,colframe=ink!45,boxrule=0.5pt,arc=1pt,left=6pt,right=6pt,breakable}

\newcommand{\blank}{\rule[-0.3em]{0pt}{1.6em}}
\newcolumntype{Y}{>{\raggedright\arraybackslash}X}
\newcommand{\steptable}[1]{\par\noindent%
  \renewcommand{\arraystretch}{1.55}%
  \begin{tabularx}{\linewidth}{@{}p{1.1cm}>{\raggedright\arraybackslash}p{7.6cm}Y@{}}
  \toprule
  단계 & 주장 & 근거 (직접 채울 것)\\ \midrule
  #1
  \bottomrule
  \end{tabularx}}
\newcommand{\Var}{\operatorname{Var}}
\newcommand{\Cov}{\operatorname{Cov}}
\newcommand{\E}{\mathbb{E}}
\newcommand{\Unif}{\operatorname{Unif}}
\newcommand{\round}{\operatorname{rd}}
\newcommand{\sX}{\sigma_X^2}
\newcommand{\sN}{\sigma_N^2}
